\documentclass[11pt,a4paper]{amsart}
\usepackage[margin=29mm]{geometry}
%\documentclass{amsart}
\usepackage{graphicx}
\usepackage[bookmarksnumbered,colorlinks,plainpages]{hyperref}
\usepackage{extarrows}
\usepackage{float}
\usepackage{amsthm}
\usepackage[all,2cell]{xy}
\usepackage{amsfonts}
\usepackage{amsmath}
\usepackage{cite}
\usepackage{amssymb}
\usepackage{xypic,leqno,amscd,amssymb,pstricks,latexsym,amsbsy,xypic,mathrsfs,verbatim}
\usepackage{multirow}
\usepackage{booktabs}
\usepackage{makecell}
\usepackage{cases}
\usepackage{appendix}
\usepackage{graphicx,colortbl}
\usepackage{tikz-cd}
\usepackage{etex}
\usepackage{mathtools}
\usepackage[bookmarksnumbered,colorlinks,plainpages]{hyperref}
\hypersetup{
	colorlinks=true,
	citecolor=blue
}
\newtheorem{theorem}{Theorem}[section]
\newtheorem{acknowledgement}[theorem]{Acknowledgement}
\newtheorem{corollary}[theorem]{Corollary}
\newtheorem{lemma}[theorem]{Lemma}
\newtheorem{proposition}[theorem]{Proposition}
\newtheorem{definition-proposition}[theorem]{Definition-Proposition}
\newtheorem{problem}[theorem]{Problem}
\newtheorem{question}[theorem]{Question}
\newtheorem{conjecture}[theorem]{Conjecture}
\theoremstyle{definition}
\newtheorem{definition}[theorem]{Definition}
\newtheorem{assumption}[theorem]{Assumption}
\newtheorem{setting}[theorem]{Setting}
\newtheorem{remark}[theorem]{Remark}
\newtheorem{example}[theorem]{Example}
\newtheorem{observation}[theorem]{Observation}
\newtheorem{construction}[theorem]{Construction}
\numberwithin{equation}{section}
\newtheorem{mainthm}{Theorem}
\renewcommand{\themainthm}{\Alph{mainthm}}

\newtheorem{innercustomthm}{{\bf Theorem}}
\newenvironment{customthm}[1]
{\renewcommand\theinnercustomthm{#1}\innercustomthm}
{\endinnercustomthm}

\newcommand{\BB}{\mathcal{B}}
\newcommand{\CC}{\mathcal{C}}
\newcommand{\OO}{{\mathcal O}}
\newcommand{\UU}{\mathcal{U}}
\newcommand{\MM}{\mathcal{M}}
\newcommand{\XX}{\mathcal{X}}
\newcommand{\PP}{\mathcal{P}}
\newcommand{\LL}{\mathcal{L}}
\newcommand{\FF}{\mathcal{F}}
\newcommand{\cs}{{\mathcal S}}
\newcommand{\Phh}{\Phi}

\newcommand{\DD}{\mathbb{D}}
\newcommand{\EE}{\mathbb{E}}
\renewcommand{\L}{\mathbb{L}}
\newcommand{\X}{\mathbb{X}}
\newcommand{\Y}{\mathbb{Y}}
\renewcommand{\P}{\mathbb{P}}
\newcommand{\Z}{\mathbb{Z}}
\newcommand{\mdots}{\begin{turn}{87}$\ddots$ \end{turn}}

\newcommand{\e}{\mathbf{e}}
\newcommand{\y}{\mathbf{y}}
\newcommand{\x}{\mathbf{x}}
\newcommand{\z}{\mathbf{z}}
\newcommand{\g}{\mathbf{g}}
\newcommand{\0}{\mathbf{0}}

\newcommand{\dimv}{\mathbf{dim}\,}
\newcommand{\ca}{\operatorname{ca}\nolimits}
\newcommand{\op}{\operatorname{op}\nolimits}
\newcommand{\Hom}{\operatorname{Hom}\nolimits}
\newcommand{\Dom}{\operatorname{Dom}\nolimits}
\newcommand{\Mod}{\operatorname{Mod}\nolimits}
\newcommand{\End}{\operatorname{End}\nolimits}
\newcommand{\Ext}{\operatorname{Ext}\nolimits}
\newcommand{\nExt}{\operatorname{nExt}\nolimits}
\newcommand{\Br}{\operatorname{Br}\nolimits}
\renewcommand{\Im}{\operatorname{Im}\nolimits}
\renewcommand{\ker}{\operatorname{Ker}\nolimits}
\newcommand{\coker}{\operatorname{Cok}\nolimits}
\newcommand{\topp}{\operatorname{top}\nolimits}
\newcommand{\soc}{\operatorname{soc}\nolimits}
\newcommand{\Cov}{\operatorname{Cov}\nolimits}
\newcommand{\bo}{\operatorname{b}\nolimits}
\newcommand{\rk}{\operatorname{rank}\nolimits}
\newcommand{\gl}{\operatorname{gl}.\operatorname{dim}\nolimits}
\newcommand{\Lotimes}{\mathop{\stackrel{\mathbf{L}}{\otimes}}\nolimits}
\newcommand{\RHom}{\mathbf{R}\strut\kern-.2em\operatorname{Hom}\nolimits}
\newcommand{\Hhom}[3]{\mathcal{H}om_{#1}(#2,#3)}
\newcommand{\sm}[1]{\left[\begin{smallmatrix} #1 \end{smallmatrix}\right]}



\DeclareMathOperator{\add}{\mathsf{add}}
\DeclareMathOperator{\ind}{\mathsf{ind}}
\DeclareMathOperator{\thick}{\mathsf{thick}}
\DeclareMathOperator{\CM}{\mathsf{CM}}
\DeclareMathOperator{\ACM}{\mathsf{ACM}}
\DeclareMathOperator{\moduleCategory}{\mathsf{mod}} \renewcommand{\mod}{\moduleCategory}
\DeclareMathOperator{\coh}{\mathsf{coh}}
\newcommand{\DDD}{\mathsf{D}}
\newcommand{\KKK}{\mathsf{K}}
\DeclareMathOperator{\qgr}{\mathsf{qgr}}
\DeclareMathOperator{\lb}{\mathsf{line}}
\DeclareMathOperator{\vect}{\mathsf{vect}}
\DeclareMathOperator{\Mon}{\mathsf{Mon}}
\DeclareMathOperator{\proj}{\mathsf{proj}}


%    Absolute value notation
\newcommand{\abs}[1]{\lvert#1\rvert}

%    Blank box placeholder for figures (to avoid requiring any
%    particular graphics capabilities for printing this document).
\newcommand{\blankbox}[2]{%
	\parbox{\columnwidth}{\centering
		%    Set fboxsep to 0 so that the actual size of the box will match the
		%    given measurements more closely.
		\setlength{\fboxsep}{0pt}%
		\fbox{\raisebox{0pt}[#2]{\hspace{#1}}}%
	}%
}



\tikzset{every picture/.style={line width=0.75pt}} %set default line width to 0.75pt   

\begin{document}
	

	\title[Cluster-roots of $\lowercase{n}$-representation-finite algebras]{Cluster-roots of $\lowercase{n}$-representation-finite algebras}

	
	\author[W. Weng] {Weikang Weng}

	\makeatletter \@namedef{subjclassname@2020}{\textup{2020} Mathematics Subject Classification} \makeatother
	
	\subjclass[2020]{16G10, 16E35}
	
	\keywords{Cluster-root, $n$-representation-finite algebra,
		$n$-cluster tilting module, Coxeter transformation,  higher Auslander algebra}

	\begin{abstract} 
	 We study cluster-roots of $n$-representation-finite algebras, namely the dimension vectors of
	 cluster-indecomposable modules. We give a direct proof of Mizuno's higher analogue of Gabriel's theorem, and show that cluster-roots are precisely the non-negative
	 integer vectors whose Coxeter orbits contain
	 the dimension vector of an indecomposable projective module.	 
	 Moreover, we give an explicit
	 description of all cluster-roots of the $(n-1)$-Auslander algebras
	 of type $A$. For $n=2$, we use periodic integer sequences to give a method for constructing counterexamples to Mizuno's positivity conjecture  within this family, which reduces to determining whether finitely many quadratic equations admit  non-negative integer solutions.
	\end{abstract}
	
	\maketitle
	
	\section{Introduction} Gabriel's theorem \cite{G} establishes a fundamental connection
	between the representation theory of finite-dimensional algebras
	and Lie theory. 
	Let $\mathbb K$ be an algebraically closed field. Gabriel's theorem states that, for a finite, connected, acyclic quiver $Q$, the
	path algebra $\mathbb KQ$ is of finite representation type if and
	only if the underlying graph of $Q$ is a Dynkin diagram of type $A$,
	$D$, or $E$. In this case, the dimension-vector map induces a
	bijection between the isomorphism classes of indecomposable
	$\mathbb KQ$-modules and the positive roots of the associated root
	system. Bernstein, Gelfand, and Ponomarev \cite{BGP} gave a
	constructive proof using reflection functors, which were later
	interpreted in terms of APR tilting \cite{APR}. 
	The connection between representations and roots was extended
	to affine quivers by Donovan--Freislich \cite{DF} and,
	independently, Nazarova \cite{N}, to valued quivers of Dynkin
	and affine type by Dlab--Ringel \cite{DR}, and to arbitrary
	finite quivers without loops by Kac \cite{Ka}.
	More recently, Geiss--Leclerc--Schr\"oer \cite{GLS} established
	a Gabriel-type theorem for algebras associated with
	symmetrizable Cartan matrices of Dynkin type, using rank
	vectors of indecomposable $\tau$-locally free modules
	in place of dimension vectors.
	Lin--Su \cite{LS} showed that, in affine type, every positive
	root is realized as the rank vector of such a module.

	Higher-dimensional Auslander--Reiten theory was introduced by Iyama
	\cite{I1} and further developed in \cite{HI1,HIO,I3,IO1,IO2,Ja}. 
	It has close connections with  several areas of mathematics, including algebraic
	geometry and singularity categories \cite{HIMO,HKV, Kv} and higher-dimensional homological algebra \cite{DKW,GKO,Ja}. 
	In higher-dimensional  Auslander--Reiten theory,  
	\emph{$n$-representation-finite algebras}, namely finite-dimensional
	algebras of global dimension at most $n$ admitting an
	$n$-cluster tilting module, are fundamental objects.
	The indecomposable direct summands of such a module are called
	\emph{cluster-indecomposable modules}.
	For $n=1$, these algebras are precisely the representation-finite
	hereditary algebras, and every indecomposable module is
	cluster-indecomposable. 
	For $n>1$, the higher Auslander algebras of linearly oriented type $A$
	form a fundamental family of examples \cite{I3}. Their recursive construction gives explicit descriptions of their quivers, relations, and higher Auslander--Reiten translations.
	Their combinatorics and connections with higher cluster categories
	have been studied in \cite{HJ,OT,W}. Dyckerhoff--Jasso--Lekili \cite{DJL} used these algebras to  study partially wrapped Fukaya categories of symmetric powers of a disk with finitely many stops. More recently, Gottesman \cite{Gott} related these algebras by derived equivalences to fractionally Calabi--Yau incidence algebras of
	lattices of order ideals of products of two chains. 

	
	Let $\Lambda$ be an $n$-representation-finite $\mathbb K$-algebra
	with acyclic quiver $Q$.
	Write $\mod\Lambda$ for the category of finite-dimensional
	right $\Lambda$-modules.
	Let $M$ be the basic $n$-cluster tilting module, and put
	$\mathcal M=\add M$, the full subcategory consisting of
	direct summands of finite direct sums of copies of $M$. 
	Denote by $\ind\mathcal M$ its full subcategory of
	indecomposable objects. Denote by $P_i$
	the indecomposable projective
	$\Lambda$-module corresponding to the vertex $i \in Q_0$. 
	Put $\ell=|Q_0|$, and write $C_\Lambda$, $q_\Lambda$ and
	$\Phi=\Phi_\Lambda$ for the Cartan matrix, Euler quadratic
	form and Coxeter transformation, respectively.
	
	Mizuno introduced analogues of the
	Bernstein--Gelfand--Ponomarev reflection functors 
	via $n$-APR tilting and used them to prove the following theorem.  
	
	
	\begin{theorem} \cite[Theorem 1.2]{Miz}\label{Mizuno thm}  The following hold. 
		\begin{itemize}
			\item[(a)]
			For each $X\in \ind\mathcal M$,  
			the dimension vector ${\dimv} X$ is a positive root of $q_\Lambda$.
			\item[(b)]
			The map $X \mapsto {\dimv} X$ gives an injection from the isomorphism classes of the objects in $\ind\mathcal M$ to the positive roots of $q_\Lambda$.
		\end{itemize}
	\end{theorem}
	
		The positive roots in the image of this map are called \emph{cluster-roots} and their set is denoted by $\mathcal R_\Lambda$.  Thus,  cluster-roots correspond bijectively to the isomorphism classes of cluster-indecomposable modules,  which gives a higher analogue of Gabriel's theorem. Indeed, when $n=1$,  cluster-roots and cluster-indecomposable modules coincide with positive roots and indecomposable modules, respectively.
		 For $n>1$, however, the dimension-vector map need not be surjective. A natural problem is therefore to characterize the cluster-roots among the positive roots of $q_\Lambda$. 
		 
		 
		 We give a direct proof of this bijection without using reflection functors, and characterize cluster-roots in terms of the Coxeter transformation as follows.  	
			
	
	
	
	\begin{theorem} [Theorem \ref{main thm}] \label{thm1}	The map $X\mapsto \dimv X$ gives a  bijection between
		the isomorphism classes of objects  in $\ind \mathcal{M}$ 
		and $\mathcal R_\Lambda$.  		
		Moreover, for any $\x\in\mathbb Z^\ell$, the following conditions are equivalent:
		\begin{enumerate}
			\item[(a)] $\x=\dimv X$ for some cluster-indecomposable $\Lambda$-module $X$;
			\item[(b)] $\x\in\mathbb Z^\ell_{\geq 0}$, $q_\Lambda(\x)=1$ and $C_{\Lambda}^{-1}\Phi^m(\x)\in\mathbb Z^\ell_{\geq 0}$ for some $m\in \Z $;		
			\item[(c)] $\x\in\mathbb Z^\ell_{\geq 0}$ and $\Phi^m(\x)=\dimv P_i$ for some $m\in \Z$ and $i\in Q_0$.
		\end{enumerate}
	\end{theorem}
	
		The proof relies on the fact that for every cluster-indecomposable module $X$, there exists a unique  pair $(t,i)\in \mathbb Z_{\geq0}\times Q_0$ such that
		$\tau_n^tX\simeq P_i$; see Proposition~\ref{cluster proper}(c). The argument then
	reduces to the projective cases, which  are uniquely determined by their
	dimension vectors; see Lemma~\ref{projective uniqueness}. 
	
	Let $d$ be the order of $\Phi$, which is finite by
	Proposition~\ref{finite order}. 
	The preceding theorem yields the following equality, thereby
	strengthening the result in
	\cite[Theorem~1.3]{Miz}.  When $n=1$, this formula recovers \cite[VII, Proposition~4.14]{ASS}.
	
	\begin{corollary} 	[Corollary \ref{computing R}]
		$
		\mathcal{R}_\Lambda
		=
		\left\{
		\Phi^m(\dimv P_i)
		\mid
		0\leq m<d,\; i\in Q_0
		\right\}
		\cap\Z_{\geq0}^{\ell}.
		$
	\end{corollary}

	
		Let $n\ge1$ and $m\ge 2$. For the higher Auslander algebra $A_m^n$ of linearly oriented
	type $A$, we use the description of 
	Oppermann--Thomas \cite{OT} with a slight modification.
	For $r\ge1$, put
	\[
	Q_0^{(r,m)}
	=
	\{(x_1,\ldots,x_r)\in\Z^r
	\mid 1\le x_1<\cdots<x_r\le m+r-1\}.
	\]
	The vertices of $A_m^n$ are indexed by $Q_0^{(n,m)}$. Write $\e_{\mathbf x}$ for the standard basis vector at a vertex
	$\mathbf x\in Q_0^{(n,m)}$.  	
	We give the following explicit description of the cluster-roots of $A_m^n$.
	
	\begin{theorem} [Theorem \ref{cluster-roots of A_m^n}] \label{thm2}
		Let $n\ge1$, $m\ge 2$ and $\Lambda=A_m^n$.
		For each $\mathbf a=(a_0,\ldots,a_n)\in Q_0^{(n+1,m)}$, put
		\[
		\rho_{\mathbf a}
		=
		\sum_{j_1=a_0}^{a_1-1}\cdots
		\sum_{j_n=a_{n-1}}^{a_n-1}
		\e_{(j_1,\ldots,j_n)}.
		\]
		Then the map $\mathbf a\mapsto\rho_{\mathbf a}$ gives a bijection
		$
		Q_0^{(n+1,m)}\rightarrow\mathcal R_\Lambda.
		$
		In particular,
		$
		|\mathcal R_\Lambda|=\binom{n+m}{n+1},
		$
		and every cluster-root has entries in $\{0,1\}$.
	\end{theorem}
	
	To characterize cluster-roots, Mizuno introduced the notion of
	$\Phi$-positivity. 
	A vector $\x$ is called \emph{$\Phi$-positive} if
	$\Phi^k(\x)\ge0$ for every $k\in\Z$.
	For even $n$, every cluster-root is $\Phi$-positive
	\cite[Theorem~5.2]{Miz}.
	For $n=2$, Mizuno conjectured the converse.
	
		\begin{conjecture}  \cite[Conjecture~1.6]{Miz}
		Let $\Lambda$ be a $2$-representation-finite algebra. Then $\Phi$-positive roots are cluster-roots.
	\end{conjecture}
	
	We show that cluster-roots are $\Phi$-positive precisely when $n$ is
	even; see Proposition \ref{positivity-comparison}. 
	For $A_m^2$, we prove that an integer vector
	$\x=(x_{ij})_{1\le i<j\le m+1}$ is $\Phi$-positive if and only if
	all $x_{ij}\ge 0$ and satisfy the triangle inequalities
	\[
	x_{ik}\le x_{ij}+x_{jk}
	\quad(1\le i<j<k\le m+1);
	\] see Proposition \ref{triangle inequalities}. 
	In particular, $L\le x_{ij}\le2L$ for some integer $L\ge1$
	implies $\Phi$-positivity of $\x$; see Corollary \ref{Squeeze}. 
	Moreover, we give a method for constructing counterexamples
	to Mizuno's conjecture by choosing
	periodic integer sequences satisfying these bounds and solving
	$q_\Lambda(\x)=1$.  This reduces to determining whether finitely many quadratic equations
	admit  non-negative integer solutions.  This construction gives the following counterexample. 
	
	
	\begin{proposition}[Proposition \ref{counterexample1}]\label{intro counterexample1}
		Let  $(g_k)_{k\in\mathbb Z}$
		be the $5$-periodic sequence determined by
		$
		(g_0,g_1,g_2,g_3,g_4)=(2,3,4,4,3).
		$ 
		Then for $\Lambda=A_{109}^2$, the vector $\x=(g_{i+j})_{1\le i<j\le110}$
		is a $\Phi$-positive root of $q_\Lambda$
		but not a cluster-root.
	\end{proposition}
	
	
	
	We now describe the organization of this paper. Section 2 recalls the Euler form, Cartan matrices, Coxeter
	transformations, and basic facts on $n$-representation-finite algebras.
	In Section 3, we give a direct proof of Mizuno's theorem, and characterize cluster-roots in terms of the Coxeter transformation, as stated in Theorem \ref{thm1}. In Section 4, we treat
	higher Auslander algebras of type $A_m^n$ and prove Theorem \ref{thm2}. In Section 5, we investigate the relationship between $\Phi$-positivity and cluster-roots. In particular, we characterize
	$\Phi$-positivity for $A_m^2$ and construct a counterexample to Mizuno's positivity conjecture.  
	
	\medskip
	\noindent\textbf{Notation.}
	Throughout this paper, we fix an integer $n\ge 1$. Let $\mathbb{K}$ be an algebraically closed field  and denote by  $D(-):=\Hom_{\mathbb{K}}(-,\mathbb{K})$.    By an algebra we mean a basic, connected and finite-dimensional $\mathbb{K}$-algebra, and by a module we mean a finite-dimensional right module.
	For an algebra $\Lambda$, we denote by $\mod \Lambda$ the category of  $\Lambda$-modules. All subcategories are assumed to be full. 		
	All quivers are assumed to be finite, connected and acyclic. For a quiver $Q$, we denote its vertex and arrow sets by $Q_0$ and $Q_1$,
	respectively, and write $\ell=|Q_0|$ for the number of vertices.
	Given a set $R$ of relations in the path algebra $\mathbb{K}Q$,
	we write $\mathbb{K}Q/(R)$ for the quotient by the ideal generated by $R$. 
	For each $i\in Q_0$, we denote by $P_i$, $I_i$ and $S_i$
	the indecomposable projective, indecomposable injective and simple
	$\mathbb{K}Q/(R)$-modules corresponding to the vertex $i \in Q_0$, respectively. 
	
	
	
	\section{Preliminaries}
		In this section, we recall some basic notions and facts needed in this paper.   		
		Throughout this paper, 	we always assume that the quiver $Q$ is finite, connected and acyclic.
	
	
	
	\subsection{Euler forms, Cartan matrices and Coxeter transformations} We recall some terminology. We refer to \cite{ASS,Miz,Rin} for more details. 	
	Let $\Lambda$ be an algebra with a complete set $\{e_i\}_{i\in Q_0}$ of primitive orthogonal idempotents. 
	Fix an integer $n\ge 1$ and assume that $\gl \Lambda\le n$. 
	
	 The \emph{Cartan matrix} of $\Lambda$ is the $(\ell\times \ell)$-matrix
	$$C_\Lambda =  (\dim_{\mathbb{K}}(e_j\Lambda e_i))_{i,j\in Q_0}.$$
		
	The \emph{Coxeter transformation} of $\Lambda$ is defined to be the matrix
	$$\Phi_\Lambda=\Phi := (-1)^nC_\Lambda^t C_\Lambda^{-1}.$$
	We  use Mizuno's sign convention \cite[Definition~2.1]{Miz}
	for the Coxeter transformation.
	This differs slightly from the usual convention; see \cite[III, Definition~3.14]{ASS}.  Since $\Lambda$ has finite global dimension, finite projective resolutions
	show that the dimension vectors of the simple modules are integral linear
	combinations of the projective dimension vectors. Hence $C_\Lambda^{-1}$
	and, consequently, $\Phi^{\pm1}$ have integral entries.
	
	The \emph{Euler form} of $\Lambda$ is the $\mathbb{Z}$-bilinear form
	$$\langle -, - \rangle : \mathbb{Z}^\ell \times \mathbb{Z}^\ell  \to \mathbb{Z}$$
	defined by $\langle \x,\y\rangle:= \x^t(C_\Lambda^{-1})^t  \y$
	for $\x ,\y  \in \mathbb{Z}^\ell $.
	We write $q_\Lambda(\x):=\langle\x,\x\rangle$ for the associated
	quadratic form. A vector $\x\in\mathbb{Z}^{\ell}$ is called
	a \emph{root} if $q_\Lambda(\x)=1$. A root $\x$ is called \emph{positive} if
	$\x\in\mathbb Z_{\geq0}^{\ell}$. 
	
	 For a $\Lambda$-module $X$, we denote the \emph{dimension vector} by 
	$${\dimv} X :=(\dim_{\mathbb{K}} (Xe_i))_{i\in Q_0}^t.$$
	Note that $C_\Lambda\mathbf e_i=\dimv P_i$ and
	$C_\Lambda^t\mathbf e_i=\dimv I_i$ for $i\in Q_0$.
	%Note that the $i$th column and row of $C_\Lambda$ are $\dimv P_i$ and $(\dimv I_i)^t$, respectively.


	\begin{lemma}\cite[Lemma 2.2]{Miz} \label{Euler form}
	Let $\Lambda=\mathbb{K}Q/(R)$ be an algebra with $\operatorname{gl.dim}\Lambda\le n$.
	\begin{itemize}
		\item[(a)]  $\Phi(\dimv P_i)=(-1)^{n}\dimv I_i$ for any $i\in Q_0$.
		\item[(b)] $\langle \x,\y\rangle=(-1)^{n}\langle \y,\Phi(\x)\rangle=\langle \Phi(\x),\Phi(\y)\rangle.$ In particular,   $q_\Lambda(\Phi(\x))=q_\Lambda(\x)$.
		\item[(c)] (Euler formula) For any $X,Y$ in $\mod\Lambda$, 
		$$ \langle {\dimv}X, {\dimv }Y\rangle  = \sum_{j \geq 0}(-1)^j\dim_{\mathbb{K}}\Ext^j_\Lambda(X,Y).$$
	\end{itemize}
	\end{lemma}
	
	\subsection{$n$-representation-finite algebras}
	The main objects of study in this paper are $n$-representation-finite
	algebras. Their theory has been developed in
	\cite{HI1,I3,IO1,IO2}, with the case $n=2$ treated in \cite{HI2}.
	
	
	
	\begin{definition}
		An algebra $\Lambda$ is called \emph{$n$-representation-finite} if $\gl\Lambda\le n$ and there exists a basic \emph{$n$-cluster tilting $\Lambda$-module $M$}, that is,
		\begin{eqnarray*}
			\add M&=&\{X\in\mod\Lambda\ |\ \Ext^i_\Lambda(M,X)=0\ \text{ for any } 0<i<n\},\\
			&=&\{X\in\mod\Lambda\ |\ \Ext^i_\Lambda(X,M)=0\ \text{ for any } 0<i<n\}.
		\end{eqnarray*}
		We call $\mathcal{M}=\add M$ the \emph{$n$-cluster tilting subcategory of $\mod \Lambda$}. An indecomposable object of  $\mathcal{M}$ is called \emph{cluster-indecomposable}, and its dimension vector is called a \emph{cluster-root}. We denote by $\ind\mathcal M$  the  subcategory of $\mathcal M$ consisting of its indecomposable objects, and set
		$$\mathcal{R}_\Lambda :=\{\dimv X \mid X \in \ind\mathcal M\}.$$
	\end{definition}
	
	Let $\Lambda$ be an algebra with $\gl\Lambda\le n$. 
	Let $\DDD^{\bo}(\mod\Lambda)$ denote the bounded derived category of $\mod\Lambda$. The \emph{Nakayama functor} is defined by\[ \nu:=D\circ\RHom_\Lambda(-,\Lambda): \DDD^{\bo}(\mod\Lambda)\rightarrow\DDD^{\bo}(\mod\Lambda).\]
	We set $\nu_n:=\nu\circ[-n]$. Since $\nu(P_i)=I_i$, the action of $\nu_n$ on the
	Grothendieck group $K_0(\DDD^{\bo}(\mod\Lambda))$ is given by $\Phi$.
	
	The \emph{$n$-Auslander--Reiten translations} are defined by the following functors: 
	\begin{align*}
		\tau_n&:=D\Ext_\Lambda^n(-,\Lambda):
		\mod\Lambda\rightarrow\mod\Lambda,\\
		\tau_n^-&:=\Ext_\Lambda^n(D\Lambda,-):
		\mod\Lambda\rightarrow\mod\Lambda.
	\end{align*}

	Let $\mathcal{M}_P$ and $\mathcal{M}_I$ be the subcategories of $\mathcal{M}$ whose objects have no non-zero projective and no non-zero injective direct summands, respectively. 	
	We collect the following important properties of $n$-representation-finite
	 algebras.  
	
	
	
	\begin{proposition}\label{cluster proper} 
		Let $\Lambda=\mathbb KQ/(R)$ be an $n$-representation-finite
		algebra, and $M$ a basic $n$-cluster tilting
		$\Lambda$-module. Set $\mathcal{M}=\add M$.  The following hold.
		\begin{itemize}
			\item[(a)] There are integers $m_i\ge0$ 
			and a  permutation $\sigma$ of $Q_0$ such that
			$\tau_n^{m_i}I_i\simeq P_{\sigma(i)}$ 	for any	$i\in Q_0$. Moreover, 
			\[
			M\simeq\bigoplus_{j\geq0}\tau_n^{-j}(\Lambda)
			\simeq\bigoplus_{j\geq0}\tau_n^j(D\Lambda).
			\] 
			\item[(b)] There are mutually quasi-inverse equivalences
			  $$\tau_n \colon \mathcal{M}_P \to \mathcal{M}_I\quad \text{and} \quad\tau_n^- \colon \mathcal{M}_I \to \mathcal{M}_P.$$ 
			\item[(c)] For each $X\in\ind\mathcal M$, there exist unique pairs
			$(t,i)$ and $(r,j)$ in $\mathbb Z_{\geq0}\times Q_0$ such that
			$\tau_n^tX\simeq P_i$ and $\tau_n^{-r}X\simeq I_j$.			
			We write $t=t(X)$ and $r=r(X)$.
			\item[(d)] $\tau_n$ and $\tau_n^-$ are nilpotent, i.e. $\tau_n^N=0$ and $\tau_n^{-N}=0$ for some integer $N\ge 1$.
			\item[(e)] 
			For $X\in\ind\mathcal M$, we have
			$\tau_nX\simeq\nu_nX$ if $X$ is non-projective, and
			$\tau_n^-X\simeq\nu_n^{-1}X$ if $X$ is non-injective.
			\item[(f)] For $X\in\ind\mathcal M$, we have
				  $$\dimv(\tau_nX)=\Phi(\dimv X) \text{ if $X$ is non-projective},  $$ 
				 and 
				 $$\dimv(\tau_n^-X)=\Phi^{-1}(\dimv X) \text{ if $X$ is non-injective}.  $$ 
				 Consequently, $\dimv(\tau_n^k X)=\Phi^k(\dimv X)$ for any
				 $0 \le k\le t(X)$, where $t(X)$ is defined in (c).
		\end{itemize}
	\end{proposition}
	\begin{proof} 
		 The statements (a), (b), (d) and (e) are given in
		\cite[Proposition~2.5]{Miz}; see also \cite{HI1,I3}.
		 The statement (c) follows from \cite[Proposition~0.2]{HI1}. 
		The first two identities in (f) are given in \cite[Lemma~3.1]{Miz}. 		
		The last identity follows by applying  the first identity repeatedly. 
	\end{proof}
	
	

	

	
	
	\section{A characterization of cluster-roots}
	In this section, we introduce the notion of $\Phi$-projective positive roots
	and give a direct proof that, for an $n$-representation-finite
	algebra, these roots coincide with the cluster-roots. 	
	Throughout this paper, let $\Lambda=\mathbb{K}Q/(R)$ be an $n$-representation-finite algebra, 
	 and  $\mathcal{M}$ the $n$-cluster tilting subcategory of $\mod\Lambda$.
	Set $\ell:=|Q_0|$, and let $\mathcal{R}_\Lambda$ denote the set of
	cluster-roots of $\Lambda$. We write $q_\Lambda$, $C_\Lambda$ and
	$\Phi=\Phi_\Lambda$ for the Euler form, Cartan matrix and Coxeter
	transformation of $\Lambda$, respectively. Denote by $\mathbf{e}_i$ the $i$th standard basis vector of $\Z^{\ell}$.
	
	

	\subsection{$\Phi$-projective positive roots and $\Phi$-orbits}	
	
	\begin{definition} \label{projective positive} A vector $\x\in\Z^\ell$ is called \emph{$\Phi$-projective} if
		$\Phi^m(\x)=\dimv P_i$ for some $m\in\Z$ and $i\in Q_0$.
		It is called \emph{$\Phi$-projective positive} if, in addition,
		$\x\in\Z_{\ge0}^\ell$.
		\end{definition}	
	
	\begin{lemma} \label{Phi-projective vector is a root}
			Each $\Phi$-projective vector is non-zero and is a root of $q_{\Lambda}$.
	\end{lemma}
	\begin{proof} Let $\mathbf{x}\in\Z^\ell$ satisfy $	\Phi^m(\x)=\dimv P_i$ for some $m\in \Z$ and $ i\in Q_0$. Then 
			\begin{alignat*}{2}
			q_\Lambda(\x)
			&=q_\Lambda(\Phi^m(\x))
			&&\qquad\text{(by Lemma~\ref{Euler form}(b))}\\
			&=q_\Lambda(\dimv P_i)
			&&\qquad\text{(by assumption)}\\
			&=\sum_{j\geq0}(-1)^j
			\dim_{\mathbb K}\Ext^j_\Lambda(P_i,P_i)
			&&\qquad\text{(by Lemma~\ref{Euler form}(c))}\\
			&=\dim_{\mathbb K}\End_\Lambda(P_i)
			&&\qquad\text{($P_i$ is projective)}\\
			&=1
			&&\qquad\text{(since $Q$ is acyclic).}
		\end{alignat*}
		Thus $\x$ is a root of $q_\Lambda$ and, in particular,
		$\x\neq0$. 
	\end{proof}
	
	The following lemma gives equivalent descriptions of $\Phi$-projectivity.
	
	\begin{lemma} \label{projectivity}
		Let $\x\in\mathbb Z^\ell$. For any $m\in\mathbb Z$, the following conditions are equivalent:
		\begin{enumerate}
			\item[(a)]  $q_\Lambda(\x)=1$ and $C_\Lambda^{-1}\Phi^m(\x)\in \Z^{\ell}_{\ge 0}$;
			\item[(b)]  $C_\Lambda^{-1}\Phi^m(\x)=\mathbf{e}_i$ for some $i\in Q_0$;
			\item[(c)] $\Phi^m(\x)=\dimv P_i$ for some  $i\in Q_0$.
		\end{enumerate}
	\end{lemma}
	
	
	
	\begin{proof}
		(a)$\Rightarrow$(b) Put
		$\g=C_\Lambda^{-1}\Phi^m(\x)$. 		
		Consider the projective module
		$
		P=\bigoplus_{i\in Q_0}P_i^{g_i}.
		$
		By the definition of the Cartan matrix, we have
		\begin{equation}\label{equ1}
			\dimv P=\sum_{i\in Q_0}g_i \dimv P_i
			=C_\Lambda \g=\Phi^m(\x)
		\end{equation}
		Therefore, we have
		\begin{alignat*}{2}
			1
			&=q_\Lambda(\x)
			&&\\
			&=q_\Lambda(\Phi^m(\x))
			&&\qquad\text{(by Lemma~\ref{Euler form}(b))}\\
			&=q_\Lambda(\dimv P)
			&&\qquad\text{(by \eqref{equ1})}\\
			&=\sum_{j\geq0}(-1)^j
			\dim_{\mathbb K}\Ext^j_\Lambda(P,P)
			&&\qquad\text{(by Lemma~\ref{Euler form}(c))}\\
			&=\dim_{\mathbb K}\End_\Lambda(P)
			&&\qquad\text{($P$ is projective)}\\
			&=\sum_{i,j\in  Q_0}
			g_i g_j\dim_{\mathbb K}\Hom_\Lambda(P_i,P_j)
			&&\\
			&\geq\sum_{i\in Q_0}g_i^2
			&&\qquad\text{(since }
			\dim_{\mathbb K}\End_\Lambda(P_i)\geq1\text{).}
		\end{alignat*}
		Since $q_\Lambda(\x)=1$, we have $\g\neq0$.
		Together with $\g\in\mathbb Z_{\geq0}^{\ell}$ and
		$\sum_{i\in Q_0}g_i^2\leq1$, this implies that
		$\g=\e_i$ for some $i\in Q_0$, proving $(b)$.
			
		(b)$\Rightarrow$(c) Multiplying $C_\Lambda^{-1}\Phi^m(\x)=\mathbf{e}_i$ by $C_\Lambda$, we obtain
		$$\Phi^m(\x)=C_\Lambda\mathbf{e}_i=\dimv P_i.$$
		
		(c)$\Rightarrow$(a) By Lemma \ref{Phi-projective vector is a root},  $q_\Lambda(\x)=1$.  Moreover, 
			$
		C_\Lambda^{-1}\Phi^m(\x)=C_\Lambda^{-1} \dimv P_i= \mathbf{e}_i\in \Z^{\ell}_{\ge 0}.
		$ 		
	\end{proof}
	
		Let
	$$\mathcal{D}:=\{\pm\dimv X\mid X\in\ind \mathcal{M}  \}.$$		
	The following proposition gives a direct proof of 	\cite[Proposition~3.5]{Miz}. 	
	\begin{proposition} \label{finite order} The finite set $\mathcal{D}$		
		is permuted by $\Phi$. If $n$ is even, the set $\mathcal{R}_{\Lambda}$ is also permuted by $\Phi$. 
		Moreover, $\Phi$ has finite order.
	\end{proposition}
	\begin{proof}
			Since $\mathcal M=\add M$, the set $\mathcal D$ is finite.
			Recall that all projective and injective $\Lambda$-modules
			belong to $\mathcal M$.
			For $X\in\ind\mathcal M$, we have
			\[
			\Phi(\dimv X)=
			\begin{cases}
				\dimv(\tau_n X), & X \text{ is non-projective},\\
				(-1)^n\dimv I_i, & X\simeq P_i,
			\end{cases}
			\]
			by Proposition~\ref{cluster proper}(f) and
			Lemma~\ref{Euler form}(a).
				It follows that
			$\Phi(\mathcal{D})\subset\mathcal{D}$.
			Since $\Phi$ is invertible and $\mathcal{D}$ is finite,
			$\Phi$ permutes $\mathcal{D}$.			
			If $n$ is even, the same formula gives
			$\Phi(\mathcal R_\Lambda)\subset\mathcal R_\Lambda$,
			so $\Phi$ also permutes $\mathcal R_\Lambda$.
			
		It remains to show the last assertion. Since $\Phi$ permutes the finite set $\mathcal D$, 			
			 there exists an integer $d\ge 1$ such that
			$\Phi^d(\x)=\x$ for any
			$\x\in\mathcal{D}$.
			In particular, $\Phi^d$ fixes each $\dimv P_i$.
			These vectors are the columns of $C_\Lambda$, so
			$\Phi^d C_\Lambda=C_\Lambda.$ 
			Since $C_\Lambda$ is invertible, $\Phi^d$ is the identity transformation.
			Hence $\Phi$ has finite order.
		\end{proof}
		
		\begin{remark} 
				Since $\Phi$ has finite order, the condition $m\in\Z$
			in the definition of $\Phi$-projectivity may equivalently
			be replaced by $m\ge 0$.
		\end{remark}
	
	



	
	
	\subsection{A characterization of cluster-roots} In this section, we give a direct proof of Mizuno's theorem and show that cluster-roots are precisely the $\Phi$-projective  positive roots. 
	
	
	
	\begin{lemma} \label{projective uniqueness} If $X\in \ind \mathcal{M}$ and $\dimv X=\dimv P_i$ for some $i\in Q_0$, then $X\simeq P_i$. 
		
	\end{lemma}
	\begin{proof}
		Since $P_i$ is projective and $Q$ is acyclic,
		by Lemma~\ref{Euler form}(c) we have
		\[
		q_\Lambda(\dimv X)
		=q_\Lambda(\dimv P_i)
		=\dim_{\mathbb K}\End_\Lambda(P_i)=1.
		\]				
		
		Assume for contradiction that $X$ is non-projective.  Put $Y=\tau_nX$.
		By Proposition~\ref{cluster proper}(f) and
		Lemma~\ref{Euler form}(a), we have
		\[
		\dimv Y=\Phi(\dimv X)
		=\Phi(\dimv P_i)=(-1)^n\dimv I_i.
		\]
		This is impossible if $n$ is odd. Hence $n$ is even
		and $\dimv Y=\dimv I_i$.
		
		For any arrow $j\to i$, the acyclicity of $Q$ implies
		that there is no path from $i$ to $j$. Then $P_ie_j=0$,
		and thus $\dimv X=\dimv P_i$ gives $Xe_j=0$.
		Similarly, for any arrow $i\to j$, there is no path
		from $j$ to $i$. Then $I_ie_j=0$, and thus
		$\dimv Y=\dimv I_i$ yields $Ye_j=0$.  		
		Thus by \cite[III, Lemma 2.2]{ASS}, we have
		$$(\topp X)e_i\simeq Xe_i \quad\text{and}\quad
		(\soc Y)e_i=Ye_i,$$ where $\topp X$ is the top of $X$ and $\soc Y$ is the socle of $Y$. The equalities
		$\dimv X=\dimv P_i$ and $\dimv Y=\dimv I_i$ imply that both  
		$Xe_i$ and $Ye_i$ are one-dimensional $\mathbb K$-spaces.
		Thus $S_i$ is a direct summand of $\topp X$ and of
		$\soc Y$. This gives a non-zero composite
		\[
		X\twoheadrightarrow \topp X \twoheadrightarrow  S_i\hookrightarrow \soc Y \hookrightarrow Y,
		\]
		 where we depict epimorphisms by $\twoheadrightarrow$ 
		and monomorphisms by $\hookrightarrow$. 
		By Serre duality and Proposition~\ref{cluster proper}(e),	we have 	
		\begin{align*}
			\Ext_\Lambda^n(X,X)&= \Hom_{\DDD^{\bo}(\mod \Lambda)}(X,X[n])\\
			&\simeq D\Hom_{\DDD^{\bo}(\mod \Lambda)}(X,\nu_nX)\\
			&\simeq D\Hom_\Lambda(X,Y)\neq0.
		\end{align*}
		Since $X\in\mathcal M$, we have
		$\Ext_\Lambda^j(X,X)=0$ for $0<j<n$.
		Since $n$ is even, the Euler formula yields
		\[
		1=q_\Lambda(\dimv X)
		=\dim_{\mathbb K}\End_\Lambda(X)
		+\dim_{\mathbb K}\Ext_\Lambda^n(X,X)
		\geq2,
		\]
		a contradiction. 		
		Thus $X$ is an indecomposable projective module.
		Since the columns of $C_\Lambda$ are linearly independent,
		$\dimv X=\dimv P_i$ implies $X\simeq P_i$.
	\end{proof}
	
	Our main result in this section is the following.  
	
	\begin{theorem} \label{main thm}
		The map $X\mapsto \dimv X$ gives a  bijection between
		the isomorphism classes of objects  in $\ind \mathcal{M}$ 
		and $\mathcal R_\Lambda$.  		
		Moreover, for any $\x\in\mathbb Z^\ell$, the following conditions are equivalent:
		\begin{enumerate}
			\item[(a)] $\x=\dimv X$ for some cluster-indecomposable $\Lambda$-module $X$;
			\item[(b)] $\x\in\mathbb Z^\ell_{\geq 0}$, $q_\Lambda(\x)=1$ and $C_{\Lambda}^{-1}\Phi^m(\x)\in\mathbb Z^\ell_{\geq 0}$ for some $m\in \Z $;		
			\item[(c)] $\x\in\mathbb Z^\ell_{\geq 0}$ and $\Phi^m(\x)=\dimv P_i$ for some $m\in \Z$ and $i\in Q_0$.
		\end{enumerate}
	\end{theorem}
	
	\begin{proof}
		We show that the map $X\mapsto \dimv X$ is a  bijection.
		Surjectivity follows from the definition of
		$\mathcal R_\Lambda$. 
		Let $X,Y\in\ind\mathcal M$ satisfy $\dimv X=\dimv Y$. Recall from  Proposition~\ref{cluster proper}(c) that for each $Z\in \ind \mathcal{M}$, there is a unique integer $t=t(Z)\geq0$
		such that $\tau_n^{t}Z$ is an indecomposable
		projective module. 
		Interchanging $X$ and $Y$ if necessary, we may assume
		that $t(X)\leq t(Y)$.
		Put $t=t(X)$ and write $\tau_n^tX\simeq P_i$.
		Since $t\leq t(Y)$, we have
		\[
		\dimv(\tau_n^tY)
		=\Phi^t(\dimv Y)
		=\Phi^t(\dimv X)
		=\dimv(\tau_n^tX)
		=\dimv P_i
		\] by Proposition~\ref{cluster proper}(f). 
		By Lemma~\ref{projective uniqueness},
		$\tau_n^tY\simeq P_i$.
		Since $\tau_n^kX$ and $\tau_n^kY$ are non-projective
		for $0\leq k<t$, applying
		Proposition~\ref{cluster proper}(b) successively yields
		\[
		X\simeq\tau_n^{-t}P_i\simeq Y.
		\]
		Thus the map is injective, and hence  bijective.
		
		It remains to prove the equivalence of (a), (b) and (c).
		
			(a)$\Rightarrow$(c) 
		Let $\x=\dimv X$ with $X\in\ind\mathcal M$.
		Then $\x\in\Z_{\geq0}^{\ell}$.
		By Proposition~\ref{cluster proper}(c), there exist
		$t\geq0$ and $i\in Q_0$ such that
		$\tau_n^tX\simeq P_i$.
		By Proposition~\ref{cluster proper}(f), (c) follows from 
		\[
		\Phi^t(\x)=\dimv(\tau_n^tX)=\dimv P_i.
		\]
		
		
		(b)$\Leftrightarrow$(c) 
		This follows directly from Lemma~\ref{projectivity}.
		
		(c)$\Rightarrow$(a) 
		Suppose that $\x\in\Z_{\geq0}^{\ell}$ and
		$\Phi^m(\x)=\dimv P_i$ for some $m\in \Z$
		and $i\in Q_0$.
		Since $\Phi$ permutes $\mathcal D$ by
		Proposition~\ref{finite order}, we have
		\[
		\x=\Phi^{-m}(\dimv P_i)\in\mathcal D.
		\]
		Hence $\x=\pm\dimv X$ for some $X\in\ind\mathcal M$.
		Since $\x\in\Z_{\geq0}^{\ell}$,
		we have $\x=\dimv X$.
	\end{proof}
	
	 
	 Let $d$ be the order of $\Phi$, which is finite by
	 Proposition~\ref{finite order}. The following corollary provides a finite procedure for computing
	 $\mathcal{R}_\Lambda$ by iterating $\Phi$ on the dimension
	 vectors of the indecomposable projective modules. A related one-sided description in terms of reflection
	 transformations is given in \cite[Theorem~4.10]{Miz}. 
	\begin{corollary} \label{computing R}
		$
		\mathcal{R}_\Lambda
		=
		\left\{
		\Phi^m(\dimv P_i) \mid
		0\leq m<d,\; i\in Q_0
		\right\}
		\cap\Z_{\geq0}^{\ell}.
		$
		\end{corollary}
	\begin{proof}
		By Theorem~\ref{main thm}, the cluster-roots are precisely the
		non-negative vectors $\Phi^{-m}(\dimv P_i)$ with $m\in \Z$ and
		$i\in Q_0$.
		Since $\Phi$ has order $d$, these powers range over
		$\Phi^0,\ldots,\Phi^{d-1}$, proving the assertion.
	\end{proof}
		
		The following example illustrates the computation of
			cluster-roots using Corollary~\ref{computing R}. 
		
	\begin{example}
		Let $\Lambda$ be the algebra presented by the  quiver 
		\[		
		1\longrightarrow2\longrightarrow3\longrightarrow4
		\] with radical square zero relations. The algebra $\Lambda$ is $3$-representation-finite.
		 Then
		\[
		C_\Lambda=
		\begin{pmatrix}
			1&0&0&0\\
			1&1&0&0\\
			0&1&1&0\\
			0&0&1&1
		\end{pmatrix}
		\quad\text{and}\quad
		\Phi
		=
		-C_\Lambda^tC_\Lambda^{-1}
		=
		\begin{pmatrix}
			0&-1&0&0\\
			0&0&-1&0\\
			0&0&0&-1\\
			1&-1&1&-1
		\end{pmatrix}.
		\]
		Moreover, $\Phi_\Lambda^5=I_4$. The dimension vectors of the
		indecomposable projective modules are
		\[
		\dimv P_1=
		\begin{pmatrix}1\\1\\0\\0\end{pmatrix},\quad
		\dimv P_2=
		\begin{pmatrix}0\\1\\1\\0\end{pmatrix},\quad
		\dimv P_3=
		\begin{pmatrix}0\\0\\1\\1\end{pmatrix},\quad
		\dimv P_4=
		\begin{pmatrix}0\\0\\0\\1\end{pmatrix}.
		\] 	
		Their $\Phi$-orbits are given by
		\[
		\begin{aligned}
			\dimv P_1
			&=
			\begin{pmatrix}1\\1\\0\\0\end{pmatrix}
			\xrightarrow{\Phi}
			\begin{pmatrix}-1\\0\\0\\0\end{pmatrix}
			\xrightarrow{\Phi}
			\begin{pmatrix}0\\0\\0\\-1\end{pmatrix}
			\xrightarrow{\Phi}
			\begin{pmatrix}0\\0\\1\\1\end{pmatrix}
			\xrightarrow{\Phi}
			\begin{pmatrix}0\\-1\\-1\\0\end{pmatrix}
			\xrightarrow{\Phi}
			\dimv P_1,
		\end{aligned}
		\]
		\[
		\begin{aligned}
			\dimv P_2
			&=
			\begin{pmatrix}0\\1\\1\\0\end{pmatrix}
			\xrightarrow{\Phi}
			\begin{pmatrix}-1\\-1\\0\\0\end{pmatrix}
			\xrightarrow{\Phi}
			\begin{pmatrix}1\\0\\0\\0\end{pmatrix}
			\xrightarrow{\Phi}
			\begin{pmatrix}0\\0\\0\\1\end{pmatrix}
			\xrightarrow{\Phi}
			\begin{pmatrix}0\\0\\-1\\-1\end{pmatrix}
			\xrightarrow{\Phi}
			\dimv P_2.
		\end{aligned}
		\]		
		These two orbits contain all the dimension vectors of the
		indecomposable projective modules.  
		Taking the non-negative vectors in these two orbits, by  Corollary~\ref{computing R}, we obtain
		\[
		\mathcal R_{\Lambda}
		=
		\left\{
		\begin{pmatrix}1\\0\\0\\0\end{pmatrix},
		\begin{pmatrix}0\\0\\0\\1\end{pmatrix},
		\begin{pmatrix}1\\1\\0\\0\end{pmatrix},
		\begin{pmatrix}0\\1\\1\\0\end{pmatrix},
		\begin{pmatrix}0\\0\\1\\1\end{pmatrix}
		\right\}.
		\]
	\end{example}
		

	
		The following result gives an explicit inverse
		of the map $X\mapsto\dimv X$
		in Theorem~\ref{main thm}.
	
	
	
	\begin{proposition}
		Let $\x\in\mathcal R_\Lambda$ and set
		\[
		s=s(\x):=
		\min\{m\in\mathbb Z_{\geq0}
		\mid C_\Lambda^{-1}\Phi^m(\x)\in \Z^{\ell}_{\ge 0}\}.
		\]
		Then there is a unique $i\in Q_0$ such that
		$
		C_\Lambda^{-1}\Phi^s(\x)=\e_i.
		$
		Under the bijection in Theorem~\ref{main thm},
		the vector $\x$ corresponds to the module
		$\tau_n^{-s}P_i$.
	\end{proposition}
	
	
	\begin{proof}
		Let $X\in\ind\mathcal M$ with $\dimv X=\x$.
		By Proposition~\ref{cluster proper}(c), there exist a unique pair $(t,j)\in \mathbb Z_{\geq0}\times Q_0$ such that 
		$\tau_n^tX\simeq P_j$.
		By Proposition~\ref{cluster proper}(f), for any $0\leq k\leq t$, 
		\[
		\dimv(\tau_n^kX)=\Phi^k(\x).
		\]
		In particular,
		$C_\Lambda^{-1}\Phi^t(\x)=\e_j$,
		so $s$ is well defined and $s\leq t$. 		
		Since $\x$ is a cluster-root, we have
		$q_\Lambda(\x)=1$. By the definition of $s$, we have 
		$C_\Lambda^{-1}\Phi^s(\x)\in\Z_{\geq0}^{\ell}$.
		Applying Lemma~\ref{projectivity} with $m=s$, we obtain
		$C_\Lambda^{-1}\Phi^s(\x)=\e_i$	 	
		for some $i\in Q_0$.
		Hence
		\[
		\dimv(\tau_n^sX)
		=\Phi^s(\x)
		=C_\Lambda\e_i
		=\dimv P_i.
		\]
		By Lemma~\ref{projective uniqueness},
		$\tau_n^sX\simeq P_i$.
		The uniqueness of $(t,j)$ gives $s=t$ and $i=j$.
	Applying $\tau_n^-$ successively to both sides of $\tau_n^sX\simeq P_i$, we have
		$
		X\simeq\tau_n^{-s}P_i.
		$
		This proves the assertion.
	\end{proof}
		
		
		
		\begin{remark}
			When $n=1$, the algebra $\Lambda$ is hereditary and
			representation-finite, and $\mathcal M=\mod\Lambda$.
			Thus $\Lambda\simeq\mathbb KQ$ for some Dynkin quiver $Q$,
			and the cluster-indecomposable modules are precisely
			the indecomposable $\Lambda$-modules. Recall from \cite[VII, Proposition 4.14]{ASS} that			
			every positive root of $q_\Lambda$ lies in the
			$\Phi$-orbit of $\dimv P_i$ for some $i\in Q_0$.
			Hence
			\[
			\mathcal R_\Lambda
			=
			\{\x\in\Z_{\geq0}^{\ell}\mid q_\Lambda(\x)=1\}.
			\]
			Thus Theorem~\ref{main thm} recovers the
			dimension-vector bijection in Gabriel's theorem.
		\end{remark}
		
	\section{Cluster-roots for higher Auslander algebras of type $A$}
		In this section, we  determine all cluster-roots of the higher
		Auslander algebras $A_m^n$.
			
		
	Let $n\geq1$ and $m\ge 2$.  Let $A_m$ be the linearly oriented
	quiver
	\[
	1\longrightarrow 2\longrightarrow\cdots\longrightarrow m,
	\]
	and set $A_m^1=\mathbb K A_m$.  For $n\geq2$, define recursively
	\[
	\Lambda=A_m^n
	=
	\operatorname{End}_{A_m^{n-1}}(M_m^{n-1}),
	\]
	where $M_m^{n-1}$ is the unique basic
	$(n-1)$-cluster tilting module of $A_m^{n-1}$.
	The algebra $A_m^n$ is called \emph{$(n-1)$-Auslander algebra of linearly oriented type $A_m$}, which  is an $n$-representation-finite algebra of global dimension $n$. For these algebras, Iyama gave explicit descriptions of their quivers with relations; see \cite{I2}.  In this paper, we follow the descriptions of $A_m^n$ in
	\cite[Theorem/Construction~3.4]{OT} with a relabelling of the
	vertices.  The algebra $A_m^n$ is presented by the quiver $Q$ with vertices
	\[
	Q_0=Q_0^{(n,m)}=
	\left\{
	\x=(x_1,\ldots,x_n)\in\mathbb Z^n
	 \mid
	1\le x_1<\cdots< x_n\le  m+n-1
	\right\},
	\]
	and arrows
	\[
	Q_1=Q_1^{(n,m)}=
	\left\{
	\x\xrightarrow{i}\sigma_i(\x)
	 \mid
	\x,\sigma_i(\x)\in Q_0,\ 1\le i\le n
	\right\},
	\]
	where $\sigma_i(\x):=(x_1,\dots,x_{i-1},x_i+1,x_{i+1},\dots,x_n)$. For
	$\x\in Q_0$ and $i,j\in\{1,\ldots,n\}$ satisfying
	$\sigma_i(\x),\sigma_j(\sigma_i(\x))\in Q_0$, the  relations are 
	\[
	(
	\x\xrightarrow{i}\sigma_i(\x)
	\xrightarrow{j}\sigma_j(\sigma_i(\x))
	)
	=
	\begin{cases}
		(
		\x\xrightarrow{j}\sigma_j(\x)
		\xrightarrow{i}\sigma_i(\sigma_j(\x))
		),
		&\text{if }\sigma_j(\x)\in Q_0,\\[2mm]
		0,&\text{otherwise}.
	\end{cases}
	\]	
	In the original labelling of \cite{OT}, the vertices are indexed by
	$n$-tuple $(a_1,\ldots,a_n)\in\mathbb Z^n$ satisfying
	\[
	1\leq a_1,\quad
	a_k+2\leq a_{k+1}\ (1\leq k<n)\quad \text{and} \quad
	a_n\leq m+2n-2.
	\] The change of variables $a_k=x_k+k-1$ identifies
	this indexing set with $Q_0$. 
	
	Clearly, the Cartan matrix $C=C_{\Lambda}=(C_{\bf u,v})_{{\bf u,v}\in Q_0}$, where 
	\begin{equation} \label{Cartan formula}
		C_{\bf u,v}
		=
		\begin{cases}
			1,
			&\text{if }v_1\le u_1<v_2\le u_2<\cdots<v_n\le u_n,\\[2mm]
			0,&\text{otherwise}.
		\end{cases}
	\end{equation}
	Indeed, $C_{\mathbf u,\mathbf v}
	=\dim_{\mathbb K}e_{\mathbf v}\Lambda e_{\mathbf u}$,
	and this space is spanned by the residue classes of paths from $\mathbf v$ to $\mathbf u$.
	The defining relations identify all such paths, whose common
	class is non-zero when
	$$v_1\le u_1<v_2\le u_2<\cdots<v_n\le u_n.$$
	Thus the space is one-dimensional in this case, and zero otherwise.
	
	\begin{example}\label{exm1}
		Let $\Lambda=A_3^2$. 
		It is presented by the following quiver, with the commutativity
		relation on the square and zero relations on small half-squares
		\[
		\begin{tikzcd}[column sep=.7em,row sep=1.2em]
			&& 14 \arrow[rd] && \\
			& 13 \arrow[ru]
			\arrow[rr,dotted,no head]
			\arrow[rd]
			&& 24 \arrow[rd] & \\
			12 \arrow[rr,dotted,no head]
			\arrow[ru]
			&& 23 \arrow[rr,dotted,no head]
			\arrow[ru]
			&& 34
		\end{tikzcd}
		\] 					
		Ordering the vertices as $(12,13,14,23,24,34)$, the Cartan matrix and Coxeter transformation are given,
	respectively, by
	\[
	C_\Lambda=
	\begin{pmatrix}
		1&0&0&0&0&0\\
		1&1&0&0&0&0\\
		1&1&1&0&0&0\\
		0&1&0&1&0&0\\
		0&1&1&1&1&0\\
		0&0&1&0&1&1
	\end{pmatrix}
	\quad\text{and}\quad
	\Phi=C_\Lambda^tC_\Lambda^{-1}=
	\begin{pmatrix}
		0&0&1&0&0&0\\
		0&0&0&0&1&0\\
		0&0&0&0&0&1\\
		1&0&-1&0&1&0\\
		0&1&-1&0&0&1\\
		0&0&0&1&-1&1
	\end{pmatrix}.
	\] By Corollary \ref{computing R}, we obtain all cluster-roots
	by iterating $\Phi$ on $\dimv P_{ij}$ for $1\le i<j\le4$,
	as follows.
		\[
		\dimv P_{12}=
		\begin{pmatrix}
			1\\1\\1\\0\\0\\0
		\end{pmatrix}
		\xrightarrow{\Phi}
		\begin{pmatrix}
			1\\0\\0\\0\\0\\0
		\end{pmatrix}
		\xrightarrow{\Phi}
		\begin{pmatrix}
			0\\0\\0\\1\\0\\0
		\end{pmatrix}
		\xrightarrow{\Phi}
		\begin{pmatrix}
			0\\0\\0\\0\\0\\1
		\end{pmatrix}
		\xrightarrow{\Phi}
		\begin{pmatrix}
			0\\0\\1\\0\\1\\1
		\end{pmatrix}
		\xrightarrow{\Phi}
		\dimv P_{12},
		\] 
		\[
		\dimv P_{13}=
		\begin{pmatrix}
			0\\1\\1\\1\\1\\0
		\end{pmatrix}
		\xrightarrow{\Phi}
		\begin{pmatrix}
			1\\1\\0\\0\\0\\0
		\end{pmatrix}
		\xrightarrow{\Phi}
		\begin{pmatrix}
			0\\0\\0\\1\\1\\0
		\end{pmatrix}
		\xrightarrow{\Phi}
		\begin{pmatrix}
			0\\1\\0\\1\\0\\0
		\end{pmatrix}
		\xrightarrow{\Phi}
		\begin{pmatrix}
			0\\0\\0\\0\\1\\1
		\end{pmatrix}
		\xrightarrow{\Phi}
		\dimv P_{13}.
		\]
		Note that we have \[
		\dimv P_{14}=
		\begin{pmatrix}
			0\\0\\1\\0\\1\\1
		\end{pmatrix},
		\quad
		\dimv P_{23}=
		\begin{pmatrix}
			0\\0\\0\\1\\1\\0
		\end{pmatrix},
		\quad
		\dimv P_{24}=
		\begin{pmatrix}
			0\\0\\0\\0\\1\\1
		\end{pmatrix},
		\quad
		\dimv P_{34}=
		\begin{pmatrix}
			0\\0\\0\\0\\0\\1
		\end{pmatrix}.
		\]
	
	\end{example}
	
	For each $\mathbf{a}=(a_0,\dots,a_{n})\in Q_0^{(n+1,m)}$, that is, $1\le a_0<\dots<a_n\le m+n$, we define
	\begin{equation*}
		\rho_{\mathbf{a}}
		:=\sum_{j_1=a_0}^{a_1-1}\cdots
		\sum_{j_n=a_{n-1}}^{a_n-1}\e_{(j_1,\ldots,j_n)}.
	\end{equation*}
	Every index $(j_1,\ldots,j_n)$ occurring in the sum belongs to
	$Q_0^{(n,m)}$, since $1\le j_1$, $j_n\le m+n-1$, and
	$j_r<a_r\le j_{r+1}$ for $1\le r<n$.
	For $\mathbf u=(u_1,\ldots,u_n)\in Q_0^{(n,m)}$, we have
	\[
	(\rho_{\mathbf a})_{\mathbf u}=
	\begin{cases}
		1,&\text{if }a_{k-1}\le u_k<a_k \
		\text{ for any }1\le k\le n,\\
		0,&\text{otherwise}.
	\end{cases}
	\]
	In particular,
	$(\rho_{\mathbf a})_{(a_0,\ldots,a_{n-1})}=1$,
	so $\rho_{\mathbf a}$ is a non-zero vector with entries in
	$\{0,1\}$.
	
	We state \cite[Theorem~3.6(1)(2)]{OT} in our notation as follows.
	
	\begin{lemma} \label{projective,injective,rho}
		Let $\x=(x_1,\ldots,x_n)\in Q_0^{(n,m)}$. Then we have
		\[
		\dimv P_{\x}=\rho_{(x_1,\ldots,x_n,m+n)}
		\quad \text{and}\quad
		\dimv I_{\x}=\rho_{(1,x_1+1,\ldots,x_n+1)}.
		\]
	\end{lemma}
	
	\begin{proof}
		We only prove the first equality, as the proof of the second is similar.
		For $\mathbf u\in Q_0^{(n,m)}$, we have 
		$(\dimv P_{\x})_{\mathbf u}=C_{\mathbf u,\x}$, which equals $1$
		precisely when
		\[
		x_1\le u_1<x_2\le u_2<\cdots<x_n\le u_n<m+n,
		\]
		and zero otherwise by \eqref{Cartan formula}. This agrees with the definition of 
		$\rho_{(x_1,\ldots,x_n,m+n)}$.
	\end{proof}
	
	\begin{proposition}  \label{Phi_rho}
		For each $\mathbf{a}=(a_0,\dots,a_{n})\in Q_0^{(n+1,m)}$, we have
		\begin{align}\label{Phi_rho formula}
			\Phi(\rho_{\mathbf{a}})
			=
			\begin{cases}
				\rho_{(a_0+1,\dots,a_n+1)},
				&\text{if } a_n<m+n,\\[2mm]
				(-1)^n\rho_{(1,a_0+1,\dots,a_{n-1}+1)},&\text{if } a_n=m+n.
			\end{cases}
		\end{align}
		Moreover, $\Phi$ has order $m+n$.
	\end{proposition}
	\begin{proof}
		If $a_n=m+n$, the preceding lemma and
		Lemma~\ref{Euler form}(a) give
		\[
		\Phi(\rho_{\mathbf a})
		=(-1)^n\dimv I_{(a_0,\ldots,a_{n-1})}
		=(-1)^n\rho_{(1,a_0+1,\ldots,a_{n-1}+1)}.
		\]		
		
		Suppose that $a_n<m+n$. First we claim that
		\begin{align*}
			\rho_{\mathbf a}
			&=\sum_{r=0}^n(-1)^{n-r}
			\dimv P_{(a_0,\ldots,\widehat{a_r},\ldots,a_n)},\\
			\rho_{(a_0+1,\ldots,a_n+1)}
			&=\sum_{r=0}^n(-1)^r
			\dimv I_{(a_0,\ldots,\widehat{a_r},\ldots,a_n)}.
		\end{align*}
		 Let
		$\mathbf u=(u_1,\ldots,u_n)\in Q_0^{(n,m)}$.
		By \eqref{Cartan formula},
		$		
		\dim_{\mathbb K}(P_{(a_0,\ldots,\widehat{a_r},\ldots,a_n)}
		e_{\mathbf u})=1
		$
		if and only if, after removing $a_r$ from $(a_0,\ldots,a_n)$,
		exactly one of the remaining $n$ entries lies in each of the intervals
		\[
		[1,u_1],\quad (u_1,u_2],\quad\ldots,\quad
		(u_{n-1},u_n].
		\] Otherwise, this dimension is zero.
		
		
		If $a_n>u_n$, the above dimension is zero for $r<n$, whereas for $r=n$ it equals $1$ precisely when
		$
		a_{k-1}\le u_k<a_k
		$ for any $1\le k\le n$. 
		Its corresponding sign is $(-1)^{n-n}=1$.	
				
		If $a_n\le u_n$ and
		the above dimension is non-zero
		for some $r$, then exactly one of the above intervals
		contains two of the entries $a_0,\ldots,a_n$, while each
		of the others contains one. The above dimension is therefore $1$ precisely when
		$r=k,k+1$ for $0\le k<n$. The corresponding contributions to the coefficient
		of $\mathbf e_{\mathbf u}$ cancel since
		\[
		(-1)^{n-k}+(-1)^{n-k-1}=0.
		\]
		Thus the first equality follows from the definition of
		$\rho_{\mathbf a}$. The second  equality follows by a similar argument.
		Applying $\Phi$ to the first equality, we obtain
		\eqref{Phi_rho formula} by Lemma~\ref{Euler form}(a)
		and the second equality. 
		
		The formula \eqref{Phi_rho formula} gives cyclic translation by $1$ modulo $m+n$,
		with a factor $(-1)^n$ whenever $m+n$ is replaced by $1$.
		After $m+n$ iterations, each entry has crossed this boundary
		once, so the total sign is $(-1)^{n(n+1)}=1$.
		Thus $\Phi^{m+n}$ fixes every $\rho_{\mathbf a}$,
		in particular all projective dimension vectors.
		These vectors are the columns of the invertible Cartan matrix
		$C_\Lambda$ and therefore form a $\mathbb Z$-basis of
		$\mathbb Z^\ell$. Thus
		 $\Phi^{m+n}=\operatorname{Id}$  is the identity transformation. Clearly, the $\Phi$-orbit of
		 $\rho_{(1,\ldots,n+1)}$ has length $m+n$ since the vectors $\rho_{\mathbf a}$, $\mathbf{a}\in Q_0^{(n+1,m)}$ are pairwise distinct. Thus $\Phi$ has order $m+n$.
	\end{proof}
	
	Now we  determine the cluster-roots of the higher
	Auslander algebras $A_m^n$.
	
	\begin{theorem} \label{cluster-roots of A_m^n}
		Let $\Lambda=A_m^n$. The map $\mathbf a\mapsto\rho_{\mathbf a}$ gives a  bijection
		$Q_0^{(n+1,m)}\to\mathcal R_\Lambda$.  In particular, we have $|\mathcal R_\Lambda|=\binom{n+m}{n+1}$ and every cluster-root has entries in $\{0,1\}$.
	\end{theorem}
	\begin{proof}
		Let $\mathbf a=(a_0,\ldots,a_n)\in Q_0^{(n+1,m)}$
		and put $s=m+n-a_n$.
		By Proposition \ref{Phi_rho} and
		Lemma~\ref{projective,injective,rho}, we have
		\[
		\Phi^s(\rho_{\mathbf a})
		=\rho_{(a_0+s,\ldots,a_{n-1}+s,m+n)}
		=\dimv P_{(a_0+s,\ldots,a_{n-1}+s)}.
		\]
		Since $\rho_{\mathbf a}$ is non-negative,
		  we have $\rho_{\mathbf a}\in\mathcal R_\Lambda$ by Theorem~\ref{main thm}.
		
		Conversely, by Theorem~\ref{main thm}, every cluster-root
		lies in the $\Phi$-orbit of an indecomposable projective
		dimension vector. By Proposition \ref{Phi_rho}, $\Phi$
		permutes the set
		\[
		\{\pm\rho_{\mathbf a}\mid
		\mathbf a\in Q_0^{(n+1,m)}\},
		\]
		which contains all projective dimension vectors.
		Thus every cluster-root equals $\pm\rho_{\mathbf a}$
		for some $\mathbf a$. Since cluster-roots and the
		vectors $\rho_{\mathbf a}$ are non-zero and non-negative,
		the sign must be positive. Hence the map is surjective.
		
		Finally, $\rho_{\mathbf a}$ uniquely determines its
		summation ranges, and hence $\mathbf a$.
		Thus the map is also injective, and hence a bijection.  In particular, $|\mathcal R_\Lambda|=\binom{n+m}{n+1}$
		since $Q_0^{(n+1,m)}$ is identified with the set of all
		 subsets of $\{1,\ldots,m+n\}$ with $(n+1)$ elements.
	\end{proof}
	
	
	
	\begin{example}
		Let $\Lambda=A_3^2$ be shown as in Example~\ref{exm1}.
		Ordering the vertices as $(12,13,14,23,24,34)$,
		the $\binom{5}{3}=10$ cluster-roots are given by the following table.
		\begingroup
		\large
		\renewcommand{\arraystretch}{1.35}
		\setlength{\arraycolsep}{8pt}
		\[
		\begin{array}{|c|c||c|c|}
			\hline
			\mathbf a & \rho_{\mathbf a}
			& \mathbf a & \rho_{\mathbf a}\\ \hline
			123 &(1,0,0,0,0,0)^{t}&145 &(0,0,1,0,1,1)^{t}\\ \hline
			124 &(1,1,0,0,0,0)^{t}&234 &(0,0,0,1,0,0)^{t}\\ \hline
			125 &(1,1,1,0,0,0)^{t}&235 &(0,0,0,1,1,0)^{t}\\ \hline
			134 &(0,1,0,1,0,0)^{t}&245 &(0,0,0,0,1,1)^{t}\\ \hline
			135 &(0,1,1,1,1,0)^{t}&345 &(0,0,0,0,0,1)^{t}\\ \hline
		\end{array}
		\]
		\endgroup
		Indeed, these follow from
		\[
		\rho_{(a,b,c)}
		=\sum_{i=a}^{b-1}\sum_{j=b}^{c-1}\mathbf e_{(i,j)}
		\quad(1\le a<b<c\le5).
		\]	
		The $\Phi$-orbits in Example~\ref{exm1} can also be computed
		via the bijection $\mathbf a\mapsto\rho_{\mathbf a}$.
		Under this bijection, $\Phi$ corresponds to adding $1$ modulo $5$
		to each index and then arranging the indices in increasing order.
		Since $\rho_{125}=\dimv P_{12}$ and
		$\rho_{135}=\dimv P_{13}$, the two orbits are
		\[
		\begin{aligned}
			\rho_{125}
			&\xrightarrow{\Phi}\rho_{123}
			\xrightarrow{\Phi}\rho_{234}
			\xrightarrow{\Phi}\rho_{345}
			\xrightarrow{\Phi}\rho_{145}
			\xrightarrow{\Phi}\rho_{125},\\
			\rho_{135}
			&\xrightarrow{\Phi}\rho_{124}
			\xrightarrow{\Phi}\rho_{235}
			\xrightarrow{\Phi}\rho_{134}
			\xrightarrow{\Phi}\rho_{245}
			\xrightarrow{\Phi}\rho_{135}.
		\end{aligned}
		\]
	\end{example}
	
	\section{$\Phi$-positivity and cluster-roots}
	In this section, we investigate the relationship between
	$\Phi$-positivity and cluster-roots.
	We show that cluster-roots are $\Phi$-positive precisely when $n$ is
	even. For $A_m^2$, we give a precise criterion for an integer vector to be
	$\Phi$-positive and give a method for constructing counterexamples to Mizuno's positivity conjecture.
	
		To characterize cluster-roots, Mizuno introduced the following notions. 
		\begin{definition} \cite[Definition 5.1]{Miz} Let $\mathbf{x}\in\Z^\ell$. 
		\begin{enumerate}
			\item[(a)] We call $\x$  \emph{$\Phi$-sign-coherent} if  $\Phi^m(\x)\in \mathbb{Z}_{\geq0}^{\ell}\cup
			\mathbb{Z}_{\leq0}^{\ell}$ for any
			$m\in\Z$.	
			\item[(b)] We call $\x$ \emph{$\Phi$-positive} if
			$\Phi^m(\x)\in\Z_{\geq0}^\ell$ for any $m\in\Z$.
		\end{enumerate}
	\end{definition}
	
	The following proposition relates cluster-roots to
	$\Phi$-sign-coherence and $\Phi$-positivity.
	
	
	
	
	\begin{proposition}\label{positivity-comparison}
		Let $\x\in\Z^\ell$.
		\begin{enumerate}
			\item[(a)]
			If $\x$ is $\Phi$-positive, then $\x$ is $\Phi$-sign-coherent.
			
			\item[(b)]
			Assume that $\x$ is a cluster-root.
			Then $\x$ is $\Phi$-sign-coherent.
			Moreover, $\x$ is $\Phi$-positive if and only if $n$ is even.
		\end{enumerate}
	\end{proposition}
	\begin{proof}
		(a) This follows from the definitions.
		
		(b) Assume that $\x$ is a cluster-root. By \cite[Theorem~5.2]{Miz}, $\x$ is
		$\Phi$-sign-coherent, and it is $\Phi$-positive
		if $n$ is even. 		
		Suppose that $n$ is odd. By Theorem~\ref{main thm},
		there are $m\in\Z$ and $i\in Q_0$ such that
		$\Phi^m(\x)=\dimv P_i$.
		By Lemma~\ref{Euler form}(a), we have
		\[
		\Phi^{m+1}(\x)
		=\Phi(\dimv P_i)
		=-\dimv I_i\notin\Z_{\geq0}^{\ell}.
		\]
		Hence $\x$ is not $\Phi$-positive.
	\end{proof}
	
	By Proposition~\ref{positivity-comparison}(b), a cluster-root is
	$\Phi$-positive if and only if $n$ is even. Thus, when $n$ is odd,
	requiring $\Phi$-positivity would exclude every cluster-root. 
	The following example shows that a $\Phi$-positive vector
	need not  be a cluster-root, even when $n$ is even.
	
	\begin{example} \label{ex:Phi-positive not cluster-root}
	We continue to discuss
	Example~\ref{exm1}. Put 
	$\x=(0,1,0,1,1,0)^t$. 
		A direct calculation gives
		$\Phi(\x)=\x$.	
		Thus $\Phi^m(\x)=\x\in\Z_{\ge 0}^{\ell}$ for any
		$m\in\mathbb Z$, so $\x$ is $\Phi$-positive.
		However, we have
		\[
		C_\Lambda^{-1}\Phi^m(\x)
		=C_\Lambda^{-1}\x=(0,1,-1,0,1,0)^t
		\notin\Z_{\geq0}^{\ell}
		\]
		for any
		$m\in\mathbb Z$. 
		Hence $\x$ is not a cluster-root by Theorem \ref{main thm}.	
	\end{example}
	
	Note that the vector $\x$ in the preceding example is not a root,
	since $q_\Lambda(\x)=2\neq 1$.
	For $\Phi$-positive roots, Mizuno posed the following conjecture.
	
	\begin{conjecture}  \cite[Conjecture 1.6]{Miz} \label{conjecture}
		Let $\Lambda$ be a $2$-representation-finite algebra. Then $\Phi$-positive roots are cluster-roots.
	\end{conjecture}
	

	
	\subsection{Triangle inequalities}
	In this subsection, we prove that $\Phi$-positivity is equivalent to
	non-negativity together with the triangle inequalities. As a direct application, we give a simple sufficient
	criterion for $\Phi$-positivity. 
	
	For the rest of this section, let $\Lambda=A_m^2$ with $m\ge2$.
	The vertex set of its quiver is  \[
	Q_0=Q_0^{(2,m)}=
	\left\{
	(i,j)
	\mid
	1\le i<j \le  m+1
	\right\}. 
	\]
	
	\begin{lemma} 
		 Let $\x=(x_{ij})_{1\le i<j\le m+1}\in \Z^{\ell}$. For any $1\le i<j \le  m+1$, we have
		 \begin{align}
		 	(C_\Lambda^{-1}\x)_{ij}
		 	&=x_{ij}-x_{i-1,j}-x_{i,j-1}+x_{i-1,j-1}, \label{C-1x}\\
		 	(\Phi\x)_{ij}&=x_{i-1,j-1}-x_{i-1,m+1}+x_{j-1,m+1},\label{Phix_ij}
		 \end{align}
		 where we use the convention that $x_{ii}=x_{0j}=0$. In particular, 
		 \begin{align}\label{q_Lambda(x)}
		 q_{\Lambda}(\x)=\sum_{1\le i <j \le m+1}(x_{ij}-x_{i+1,j})(x_{ij}-x_{i,j-1}).
		 \end{align}
		
	\end{lemma}
	\begin{proof} Put $\g=C_\Lambda^{-1}\x$, so that $\x=C_\Lambda\g$.
		By \eqref{Cartan formula},  we have
		\[
		x_{ij}
		=\sum_{1\le a<b\le m+1}C_{ij,ab}g_{ab}
		=\sum_{a=1}^{i}\sum_{b=i+1}^{j}g_{ab},
		\]
		since $C_{ij,ab}$ is $1$ if $a\le i<b\le j$,
		and $0$ otherwise.
		Hence
		\[
		x_{ij}-x_{i,j-1}=\sum_{a=1}^{i}g_{aj}
		\quad\text{and}\quad
		x_{i-1,j}-x_{i-1,j-1}=\sum_{a=1}^{i-1}g_{aj}.
		\]
		Thus we have
		$(C_\Lambda^{-1}\x)_{ij}
		=g_{ij}
		=x_{ij}-x_{i-1,j}-x_{i,j-1}+x_{i-1,j-1}$. 		
		Since $\Phi=C_\Lambda^tC_\Lambda^{-1}$, we obtain
		\[
		(\Phi\x)_{ij}
		=\sum_{a=i}^{j-1}\sum_{b=j}^{m+1}
		g_{ab}
		=x_{i-1,j-1}+x_{j-1,m+1}-x_{i-1,m+1}.
		\]
		Finally, substituting \eqref{C-1x} into
		$q_\Lambda(\x)=\x^t (C_\Lambda^{-1})^t\x=\x^tC_\Lambda^{-1}\x$, we obtain
		\[
		\begin{aligned}
			q_\Lambda(\x)
			&=\sum_{1\le i<j\le m+1}
			x_{ij}\bigl(x_{ij}-x_{i-1,j}
			-x_{i,j-1}+x_{i-1,j-1}\bigr)\\
			&=\sum_{1\le i<j\le m+1}
			x_{ij}(x_{ij}-x_{i,j-1})
			-\sum_{1\le i<j\le m+1}
			x_{ij}(x_{i-1,j}-x_{i-1,j-1})\\
			&=\sum_{1\le i<j\le m+1}
			x_{ij}(x_{ij}-x_{i,j-1})
			-\sum_{1\le i<j\le m+1}
			x_{i+1,j}(x_{ij}-x_{i,j-1})\\
			&=\sum_{1\le i<j\le m+1}
			(x_{ij}-x_{i+1,j})(x_{ij}-x_{i,j-1}).
		\end{aligned}
		\]
		In the third equality, we replace $i$ by $i+1$ in the second
		sum, giving the range $0\le i<j-1$ for each fixed $j$.
		We then restore the original range $1\le i<j$ by omitting the term
		with $i=0$ and adding the term with $i=j-1$. These terms vanish
		since $x_{0j}=x_{0,j-1}=0$ and $x_{jj}=0$, respectively.
	\end{proof}
	
	The following result gives a precise criterion for an integer vector to be
	$\Phi$-positive. 
	
	\begin{proposition}\label{triangle inequalities}
		Let $\x=(x_{ij})_{1\le i<j\le m+1}\in \Z^{\ell}$. Then $\x$ is $\Phi$-positive if and only if $\x\in \Z^{\ell}_{\ge 0}$, and for any $1\le i<j<k\le m+1$,
		$$x_{ik}\le x_{ij}+x_{jk}.$$
	\end{proposition}
	\begin{proof}
		Suppose first that $\x$ is $\Phi$-positive. Then $\x\in \Z^{\ell}_{\ge 0}$.
		For $1\le i<j<k\le m+1$, set $x_{ii}=0=x_{0j}$ and $r=m+2-k$. We show  that
		\begin{align} \label{Phi^r}
			(\Phi^r\x)_{i+r,j+r}=x_{ij}+x_{jk}-x_{ik}
		\end{align}
		by induction on $r$. 
		The case $r=1$ follows directly from \eqref{Phix_ij}.
		For $r>1$, applying \eqref{Phix_ij} and the induction hypothesis
		to $(i,j,k+1)$, $(j,k,k+1)$ and $(i,k,k+1)$ gives
		\[
		\begin{aligned}
			(\Phi^r\x)_{i+r,j+r}
			&=(x_{ij}+x_{j,k+1}-x_{i,k+1})
			+(x_{jk}+x_{k,k+1}-x_{j,k+1})
			-(x_{ik}+x_{k,k+1}-x_{i,k+1})\\
			&=x_{ij}+x_{jk}-x_{ik}.
		\end{aligned}
		\]
		This completes the induction.  Since $\x$ is $\Phi$-positive,  we have $x_{ij}+x_{jk}-x_{ik}\ge0$ by \eqref{Phi^r}. 
		
		Conversely, suppose that $\x\in \Z^{\ell}_{\ge 0}$ and
		$x_{ik}\le x_{ij}+x_{jk}$ for all $1\le i<j<k\le m+1$.
		Put $\mathbf y=\Phi\x$. By \eqref{Phix_ij}, we have
		$\mathbf y\in \Z^{\ell}_{\ge 0}$ and
		\[
		y_{ij}+y_{jk}-y_{ik}
		=
		\begin{cases}
			x_{j-1,k-1},&i=1,\\
			x_{i-1,j-1}+x_{j-1,k-1}-x_{i-1,k-1},&i>1.
		\end{cases}
		\]
		Thus $\mathbf y$ satisfies the same inequalities.
		It follows by induction that $\Phi^h\x\in \Z^{\ell}_{\ge 0}$ for all $h\ge0$.
		Since $\Phi^{m+2}=\operatorname{Id}$ by
		Proposition~\ref{Phi_rho}, this holds for every $h\in\mathbb Z$.
		Hence $\x$ is $\Phi$-positive.	
	\end{proof}
	
	As a direct application of  Proposition \ref{triangle inequalities}, we obtain the following simple sufficient
	criterion for $\Phi$-positivity. 
	
		\begin{corollary} \label{Squeeze}
		Let $\x=(x_{ij})_{1\le i<j\le m+1}\in\Z^\ell$. 
		If there exists an integer $L\ge1$ such that
		\[
		L\le x_{ij}\le2L
		\] holds for any $1\le i<j\le m+1$,
		then $\x$ is $\Phi$-positive.
	\end{corollary}
	\begin{proof}
		Since $L\ge1$, we have $\x\in\Z_{\ge0}^{\ell}$. Moreover,
		for any $1\le i<j<k\le m+1$,
		\[
		x_{ik}\le2L\le x_{ij}+x_{jk}.
		\]
		Thus $\x$ is $\Phi$-positive by Proposition \ref{triangle inequalities}.
	\end{proof}
	
	
	
	
	
	\subsection{A counterexample to Mizuno's positivity conjecture} 
	In this subsection,  we use
	periodic integer sequences to construct a counterexample to Mizuno's
	positivity conjecture.
	
	Let $(g_k)_{k\in\mathbb Z}$ be a
	$p$-periodic sequence of integers with $p\ge1$, that is,
	$g_{k+p}=g_k$ for all $k\in\mathbb Z$.
	For $m\ge2$, put $\Lambda=A_m^2$ and set
	\[
	x_{ij}=g_{i+j}
	\quad (1\le i<j\le m+1).
	\]

	
	As before, we use the convention that $x_{ii}=x_{0j}=0$. By \eqref{q_Lambda(x)}, we have	
	\begin{align}
		q_\Lambda(\x)
		&=\sum_{1\le i<j\le m+1}  \label{q_x periodic}
		(x_{ij}-x_{i+1,j})(x_{ij}-x_{i,j-1})\\
		&=\sum_{i=1}^{m}x_{i,i+1}^{\,2} \nonumber
		+\sum_{\substack{1\le i<j\le m+1\\j-i\ge2}}
		(x_{ij}-x_{i+1,j})(x_{ij}-x_{i,j-1})\\
		&=\sum_{i=1}^{m}g_{2i+1}^{\,2} \nonumber
		-\sum_{\substack{1\le i<j\le m+1\\j-i\ge2}}
		(g_{i+j+1}-g_{i+j})(g_{i+j}-g_{i+j-1}). 
	\end{align}
	
	
		For each integer $r\ge2$, let $J_r$ be the last expression in \eqref{q_x periodic} with $m$ replaced by $r$.  Set $J_0:=0$ and $J_1:=g_3^2$.
		Thus $J_m=q_\Lambda(\x)$. We set
		$$ H:=\sum_{k=0}^{p-1}(g_{k+1}-g_k)(g_k-g_{k-1}).$$

	\begin{lemma}\label{periodic recurrence}
		For any integer $r\ge0$, we have
		$J_{r+p}=J_r+J_p-rH.$
		Consequently, if $m=bp+a$ with $0\le a<p$ and $b\ge0$, then
		\begin{equation}\label{periodic residue formula}
		J_m
		=
		-\frac{pH}{2}b^2
		+\left(\frac{p-2a}{2}H+J_p\right)b
		+J_a.
		\end{equation}
	\end{lemma}
	
	\begin{proof}
		A straightforward calculation yields 
		\begin{align}	\label{S_r}	
		S_r:=J_{r+1}-J_r
		=
		g_{2r+3}^{\,2}
		-\sum_{k=r+3}^{2r+2}(g_{k+1}-g_k)(g_k-g_{k-1}),
	\end{align}
		Periodicity gives $S_{r+p}=S_r-H$.
		Hence
		$$
		\begin{aligned}
		J_{r+p}-J_r
		&=(J_p+\sum_{j=0}^{r-1}S_{j+p})-\sum_{j=0}^{r-1}S_j\\
		&=J_p+\sum_{j=0}^{r-1}(S_{j+p}-S_j)\\
		&=J_p-rH.
		\end{aligned}
		$$
		 Iterating the preceding recurrence, we obtain  
		 \[
		 \begin{aligned}
		 J_m=	J_{bp+a}
		 	&=J_a+\sum_{k=0}^{b-1}
		 	\bigl(J_p-(a+k p)H\bigr)\\
		 	&=J_a+b(J_p-aH)-pH\binom{b}{2}\\
		 	&=-\frac{pH}{2}b^2
		 	+\left(\frac{p-2a}{2}H+J_p\right)b+J_a.
		 \end{aligned}
		 \] Thus we have the assertion.
	\end{proof}


	
		To construct counterexamples, we assume that the chosen $p$-periodic sequence $(g_k)_{k\in\mathbb Z}$ satisfies $H>0$ and that there exists an
		integer $L\ge2$ such that
		\[
		L\le g_k\le2L
		\quad\text{for all }k\in\mathbb Z.
		\] 
	Then $\x$ is $\Phi$-positive by Corollary~\ref{Squeeze}.
	Since every coordinate of $\x$ is at least $2$,
	it is not a cluster-root by
	Theorem \ref{cluster-roots of A_m^n}.
	Thus every solution of $J_m=1$ with $m\ge2$ gives a counterexample to Conjecture~\ref{conjecture}.
	
	By Lemma~\ref{periodic recurrence}, it suffices to compute
	$J_1,\dots,J_p$.
	For each $a\in\{0,\dots,p-1\}$, set
	\[
	A=pH,\quad
	B_a=(p-2a)H+2J_p\quad\text{and}\quad
	\Delta_a=B_a^2+8A(J_a-1).
	\]
	Equation \eqref{periodic residue formula} gives
	\[
	J_{bp+a}=1
	\iff
	Ab^2-B_ab+2(1-J_a)=0.
	\]
	Since $A>0$, a non-negative integer solution exists if and only if
	$\Delta_a=c^2$ for some integer $c\ge0$ and at least one of
	$
	b=(B_a\pm c)/{2A}
	$ 
	is a non-negative integer.
	
	
	We may further restrict to sequences satisfying
	$\gcd(g_0,\dots,g_{p-1})=1$, since otherwise every $J_m$
	is divisible by the square of this common divisor.
	
	\begin{remark}
		The condition $H>0$ cannot hold if  $p\le4$. Hence these cases cannot produce counterexamples by the
		above construction. Indeed, 	
		let
		$
		\delta_k=g_k-g_{k-1}.
		$
		Then
		$
		H=\sum_{k=0}^{p-1}\delta_{k+1}\delta_k.
		$			The cases $p=1$ and $p=2$ are included in the case $p=4$,
		since every $1$- or $2$-periodic sequence is also $4$-periodic. 		
		Assume  $p=3$. Since
		$\delta_0+\delta_1+\delta_2=0$, we obtain
		\[
		H
		=\delta_0\delta_1+\delta_1\delta_2+\delta_2\delta_0
		=-\frac12\left(\delta_0^2+\delta_1^2+\delta_2^2\right)
		\le0.
		\] 
				
		Finally, assume $p=4$. Using
		$\delta_0+\delta_1+\delta_2+\delta_3=0$, we have
		\[
		\begin{aligned}
			H
			=\delta_0\delta_1+\delta_1\delta_2
			+\delta_2\delta_3+\delta_3\delta_0
			=(\delta_0+\delta_2)(\delta_1+\delta_3)
			=-(\delta_0+\delta_2)^2\le0.
		\end{aligned}
		\]
	\end{remark}
	
		Now we give an explicit $5$-periodic sequence
		that produces a counterexample.
	
	\begin{proposition}\label{counterexample1}
		Choose $p=5$ and $L=2$, and define the $5$-periodic sequence
		$(g_k)_{k\in\mathbb Z}$ by
		\[
		(g_0,g_1,g_2,g_3,g_4)=(2,3,4,4,3) \quad \text{and}
		\quad g_{k+5}=g_k.
		\] Then for $\Lambda=A_{109}^2$ the vector 
		$\x=(g_{i+j})_{1\le i<j\le110}$ is a $\Phi$-positive root but not a cluster-root. 
		\end{proposition}
		\begin{proof}
	 The values of $(g_{k+1}-g_k)(g_k-g_{k-1})$ for $k=0,\dots,4$
				are $-1,1,0,0,1$, respectively. Hence $H=1$.
		By definition of $J_r$ $(r\ge 0)$ and \eqref{S_r}, we have
		\[
		\begin{aligned}
			J_0&=0, &\quad J_1&=16,\\
			J_2&=16+4-1=19, &\quad
			J_3&=19+16=35,\\
			J_4&=35+9-1=43, &\quad
			J_5&=43+9=52.
		\end{aligned}
		\]
		Hence $A=5$, $B_a=109-2a$ and $$\Delta_a=(109-2a)^2+40(J_a-1).$$ 		
			Among $\Delta_0,\dots,\Delta_4$, one can check that only $\Delta_4$ is a perfect square, namely $\Delta_4=109^2$.
		The corresponding values of $b$ are
		\[
		b=\frac{101\pm109}{10}=21\text{ or }-\frac45.
		\]
		Thus $b=21$ and $m=5b+4=109$. In particular, by Lemma \ref{periodic recurrence}, we have
		\[
			J_{109}
			=-\frac{5}{2}\cdot 21^2
			+\left(\frac{5-8}{2}+52\right)\cdot 21+43
			=1.
		\]
		Hence $q_\Lambda(\x)=1$. 
		Since $2\le x_{ij}\le4$ for all $1\le i<j\le110$,
		the vector $\x$ is $\Phi$-positive by Corollary~\ref{Squeeze}. 
		On the other hand, every coordinate of $\x$ is at least $2$,
		whereas every cluster-root of $A_{109}^2$ has entries in
		$\{0,1\}$ by Theorem~\ref{cluster-roots of A_m^n}.
		Therefore, $\x$ is not a cluster-root.
	\end{proof}

	
	

	


		
		\begin{thebibliography}{99}
		\bibitem{ASS}
		I.~Assem, D.~Simson and A.~Skowro\'nski,
		\emph{Elements of the Representation Theory of Associative Algebras.
			Vol.~1: Techniques of Representation Theory},
		London Mathematical Society Student Texts, vol.~65,
		Cambridge University Press, Cambridge, 2006.
		
		\bibitem{APR}
		M.~Auslander, M.~I.~Platzeck and I.~Reiten,
		\emph{Coxeter functors without diagrams},
		Trans. Amer. Math. Soc. {250} (1979), 1--46.
		
		\bibitem{BGP}
		I.~N.~Bernstein, I.~M.~Gelfand and V.~A.~Ponomarev,
		\emph{Coxeter functors and Gabriel's theorem},
		Russian Math. Surveys {28} (1973), 17--32.
		
		
		\bibitem{DR}
		V.~Dlab and C.~M.~Ringel,
		\emph{Indecomposable representations of graphs and algebras},
		Mem. Amer. Math. Soc. {6} (173) (1976).
		
		\bibitem{DF}
		P.~Donovan and M.~R.~Freislich,
		\emph{The representation theory of finite graphs and associated algebras},
		Carleton Mathematical Lecture Notes 5,
		Carleton University, Ottawa, 1973.
		
		
		\bibitem{DJL}
		T.~Dyckerhoff, G.~Jasso and Y.~Lekili,
		\emph{The symplectic geometry of higher Auslander algebras:
			symmetric products of disks},
		Forum Math. Sigma {9} (2021), e10.
		
		\bibitem{DKW}
		T.~Dyckerhoff, G.~Jasso and T.~Walde,
		\emph{Simplicial structures in higher Auslander--Reiten theory},
		Adv. Math. {355} (2019), 106762.
		
		\bibitem{G}
		P.~Gabriel,
		\emph{Unzerlegbare Darstellungen I},
		Manuscripta Math. {6}  (1972), 71--103.
		
		\bibitem{GKO}
		C.~Geiss, B.~Keller and S.~Oppermann,
		\emph{$n$-angulated categories},
		J. Reine Angew. Math. {675} (2013), 101--120.
		
		\bibitem{GLS}
		C.~Geiss, B.~Leclerc and J.~Schr\"oer,
		\emph{Quivers with relations for symmetrizable Cartan matrices
			I: Foundations},
		Invent. Math. {209}  (2017), 61--158.
		
		\bibitem{HI1}
		M.~Herschend and O.~Iyama,
		\emph{$n$-representation-finite algebras and twisted fractionally
			Calabi--Yau algebras},
		Bull. Lond. Math. Soc. {43}  (2011), 449--466.
		
		\bibitem{HI2}
		M.~Herschend and O.~Iyama,
		\emph{Selfinjective quivers with potential and
			$2$-representation-finite algebras},
		Compos. Math. {147} (2011), 1885--1920.
		
		\bibitem{HIMO}
		M.~Herschend, O.~Iyama, H.~Minamoto and S.~Oppermann,
		\emph{Representation theory of Geigle--Lenzing complete intersections},
		Mem. Amer. Math. Soc. {285} (2023), no.~1412.
		
		\bibitem{HIO}
		M.~Herschend, O.~Iyama and S.~Oppermann,
		\emph{$n$-representation infinite algebras},
		Adv. Math. {252} (2014), 292--342.
		
		\bibitem{HJ}
		M.~Herschend and P.~J\o rgensen,
		\emph{Classification of higher wide subcategories for higher Auslander
			algebras of type $A$},
		J. Pure Appl. Algebra {225}  (2021), 106583.
		
		\bibitem{HKV}
		M.~Herschend, S.~Kvamme and L.~Vaso,
		\emph{$n\mathbb Z$-cluster tilting subcategories for Nakayama algebras},
		Math. Z. {309} (2025), no.~37.
		
		\bibitem{I2}
		O.~Iyama,
		\emph{Auslander correspondence},
		Adv. Math. {210}  (2007), 51--82.
		
		\bibitem{I1}
		O.~Iyama,
		\emph{Higher-dimensional Auslander--Reiten theory on maximal
			orthogonal subcategories},
		Adv. Math. {210}  (2007), 22--50.
		
		\bibitem{I3}
		O.~Iyama,
		\emph{Cluster tilting for higher Auslander algebras},
		Adv. Math. {226}  (2011), 1--61.
		
		\bibitem{IO1}
		O.~Iyama and S.~Oppermann,
		\emph{$n$-representation-finite algebras and $n$-APR tilting},
		Trans. Amer. Math. Soc. {363}  (2011), 6575--6614.
		
		\bibitem{IO2}
		O.~Iyama and S.~Oppermann,
		\emph{Stable categories of higher preprojective algebras},
		Adv. Math. {244} (2013), 23--68.
		
		\bibitem{Ja}
		G.~Jasso,
		\emph{$n$-abelian and $n$-exact categories},
		Math. Z. {283} (2016), 703--759.
		
		\bibitem{Gott}
		T.~Gottesman,
		\emph{Fractionally Calabi--Yau lattices that tilt to higher
			Auslander algebras of type $A$},
		Adv. Math. \textbf{488} (2026), 110785.
		
		\bibitem{Ka}
		V.~G.~Kac,
		\emph{Infinite root systems, representations of graphs and invariant theory},
		Invent. Math. {56}  (1980), 57--92.
		
		\bibitem{Kv}
		S.~Kvamme,
		\emph{$d\mathbb Z$-cluster tilting subcategories of singularity categories},
		Math. Z. {297}  (2021), 803--825.
		
		\bibitem{LS}
		Z.~Lin and X.~Su,
		\emph{Affine root systems, stable tubes, and a conjecture by
			Geiss--Leclerc--Schr\"oer},
		Int. Math. Res. Not. IMRN {2025}  (2025), rnae279.
		
		\bibitem{Miz}
		Y.~Mizuno,
		\emph{A Gabriel-type theorem for cluster tilting},
		Proc. Lond. Math. Soc. (3) {108}  (2014), 836--868.
		
		\bibitem{N}
		L.~A.~Nazarova,
		\emph{Representations of quivers of infinite type},
		Izv. Akad. Nauk SSSR Ser. Mat. {37}  (1973), 752--791.
		
		\bibitem{OT}
		S.~Oppermann and H.~Thomas,
		\emph{Higher-dimensional cluster combinatorics and representation theory},
		J. Eur. Math. Soc. {14}  (2012), 1679--1737.
		
		\bibitem{Rin}
		C.~M.~Ringel,
		\emph{Tame Algebras and Integral Quadratic Forms},
		Lecture Notes in Mathematics, vol.~1099,
		Springer-Verlag, Berlin, 1984.
		
		\bibitem{W}
		N.~J.~Williams,
		\emph{New interpretations of the higher Stasheff--Tamari orders},
		Adv. Math. {407} (2022), 108552.
		\end{thebibliography}
		
	
		
		
			\vskip 5pt
		\noindent {\scriptsize   \noindent  Weikang Weng\\
			School of Mathematical Sciences, \\
			Xiamen University, Xiamen, 361005, Fujian, China.\\
			E-mail: 
			wkweng@stu.xmu.edu.cn\\ }
		\vskip 3pt
		
	\end{document}
	
	
	
	
	
	%------------------------------------------------------------------------------
	% End of journal.tex
	%------------------------------------------------------------------------------
